\title{Direct Preference Optimization: Your Language Model is Secretly a Reward Model}

\begin{abstract}
Large-scale unsupervised language models (LMs) acquire extensive world knowledge and demonstrate some capability for reasoning, yet guiding their outputs with precision remains challenging due to the absence of supervision during training. Current strategies to increase controllability usually involve gathering human assessments of generation quality and adapting the LM to favor preferred outputs, commonly through reinforcement learning from human feedback (RLHF). Nevertheless, RLHF presents significant complexity and instability: it requires first training a reward model that approximates human judgments, followed by applying reinforcement learning to adjust the original LM for higher estimated reward, all while preventing substantial divergence from the pretrained distribution. This paper introduces a novel reward model formulation within RLHF, allowing the associated optimal policy to be expressed in closed form. This reparameterization reduces the RLHF optimization to a straightforward classification objective. Our method, termed \textit{Direct Preference Optimization} (DPO), offers a stable, effective, and resource-efficient solution, removing the necessity for language model sampling during adaptation and obviating extensive hyperparameter searches. Empirical results demonstrate that DPO matches or outperforms existing preference alignment approaches for LMs, including PPO-based RLHF. Specifically, DPO surpasses PPO-based RLHF in modulating the sentiment of generated text and achieves equivalent or superior results in tasks such as summarization and single-turn dialogue, while offering greater simplicity in implementation and training.
\end{abstract}

\section{Introduction}
Large-scale unsupervised language models (LMs), when trained on expansive corpora, exhibit a range of remarkable capabilities~\citep{chowdhery2022palm, brown2020language, touvron2023llama,bubeck2023sparks}. Despite their impressive performance, these models draw on data created by humans with diverse motivations, expertise, and intentions. Not all such human-derived behaviors are appropriate for our purposes; for example, in an AI-based code assistant, it is beneficial for the model to \textit{recognize} typical programming errors for the sake of correction, yet we prefer that it emulates the rare, high-caliber coding proficiencies reflected in only a subset of its training data when actually generating code. Similarly, while a language model should be able to \textit{identify} prevalent misconceptions—such as one believed by half of individuals—it is undesirable for the model to echo this misconception as fact in 50\% of its answers regarding the topic. Thus, distinguishing and prioritizing a model’s \emph{intended outputs and behaviors} from the expansive spectrum of its \textit{knowledge and skills} is essential to crafting AI systems that are safe, effective, and amenable to human control \citep{ouyang2022training}. Although present techniques generally employ reinforcement learning (RL) to align LMs with human preferences, this work will demonstrate that the RL objective frequently utilized in current practice is, in fact, precisely optimized via a straightforward binary cross-entropy objective, thereby streamlining the entire preference learning process.

Broadly, prevailing strategies impart preferred behaviors to language models through the use of well-curated datasets reflecting human preferences regarding safety and usefulness. This phase of preference learning takes place after a preliminary period of large-scale, unsupervised pre-training over massive text datasets. The simplest preference learning paradigm involves supervised fine-tuning on high-quality responses generated by humans. Nonetheless, the most successful methodology to date remains reinforcement learning from human (or AI) feedback (RLHF/RLAIF; \citep{christiano2017deep,bai2022constitutional}). RLHF leverages a reward model, fit to a collection of human preference data, and subsequently applies RL to tune a language model’s policy so that it generates responses with higher reward, all while avoiding excessive deviation from the pre-trained model. While RLHF can endow models with outstanding conversational and programming skills, the associated pipeline is significantly more intricate than supervised learning, necessitating the training of multiple language models and in-loop sampling from the policy, resulting in considerable computational demands.

This work demonstrates a method for optimizing language models to align with human preferences, without resorting to explicit reward modeling or reinforcement learning techniques. We introduce \textit{Direct Preference Optimization} (DPO), an approach that implicitly targets the same objective as prevalent RLHF algorithms—namely, maximizing reward subject to a KL-divergence constraint—while remaining both easy to implement and train. In essence, DPO adjusts the relative log likelihood between preferred and non-preferred outputs by applying a dynamically determined, example-specific importance weight. This adjustment mitigates the model collapse observed when applying a naive probability ratio objective. Similar to prior methods, DPO utilizes a theoretical framework for preference modeling, such as the Bradley-Terry model \cite{bradley1952rankanalysis}, to evaluate how closely a reward function matches observed human preference data. Distinct from established approaches that use preference models to create a loss function for training a reward model, followed by separate policy optimization with respect to this learned reward, DPO leverages a change of variables to express the preference loss directly in terms of the policy parameters. Leveraging a corpus of human preference data over generated responses, DPO enables direct policy training via a straightforward binary cross entropy loss, resulting in a policy that is optimal for an implicitly defined reward derived from the preference annotations.

Our principal contribution is the introduction of Direct Preference Optimization (DPO), a straightforward, reinforcement learning-free technique for optimizing language models according to preference data. Experimental results indicate that DPO matches or exceeds the performance of contemporary methods, including those based on PPO-style RLHF, on preference-driven tasks such as sentiment control, summarization, and dialogue generation, across language models with up to 6B parameters.

Self-supervised language models that have grown in scale exhibit the ability to solve certain tasks in a zero-shot fashion \citep{radford2019language}, or when provided with only a handful of prompts (few-shot learning) \citep{gpt3,megatron,chowdhery2022palm}. Nevertheless, refining their abilities on downstream applications and achieving better alignment with user expectations often requires additional fine-tuning on corpora composed of instructions accompanied by human-generated completions \citep{mishra-etal-2022-cross,sanh2022multitask,chung2022scaling,thoppilan2022lamda}. This process, known as 'instruction-tuning,' not only increases the flexibility of LLMs to handle instructions absent from the fine-tuning dataset, but also tends to enhance overall practical applicability \citep{chung2022scaling}. In spite of the progress made through instruction tuning, gathering \textit{relative} assessments from humans regarding the quality of model responses proves to be less demanding than collecting high-quality expert demonstrations. Accordingly, recent research has pursued further fine-tuning using datasets that encode human preference information, leading to improvements across domains such as translation \citep{kreutzer-etal-2018-reliability}, summarization \citep{stiennon2022learning,ziegler2020finetuning}, narrative generation \citep{ziegler2020finetuning}, and instruction adherence \citep{ouyang2022training,ramamurthy2023is}. These approaches typically begin by learning a neural reward model aligned with the observed preference data, frequently operationalized using models like Bradley-Terry \citep{bradley1952rankanalysis}. Following this, the language model is adjusted to maximize the reward through reinforcement learning techniques, including REINFORCE \citep{williams1992reinforce}, proximal policy optimization (PPO; \cite{schulman2017proximal}), or related algorithms \citep{ramamurthy2023is}. Another related research avenue employs LLMs previously tuned via instruction following and augmented with human feedback to generate synthetic datasets reflecting preferences for qualities such as safety or innocuousness \citep{bai2022constitutional}, utilizing minimal direct human oversight—often limited to providing rubric-style text guidance for the model's own labeling. Collectively, these strategies represent the intersection between two distinct threads: the use of reinforcement learning to train language models on a variety of goals~\citep{Ranzato2015SequenceLT,paulus2018a,wu2018learning}, and the design of frameworks that enable learning from human preference signals \citep{christiano2017deep,kupcsik2018learning}. While harnessing relative human preferences is intuitively attractive, the process of optimizing large language models via reinforcement learning introduces substantial operational difficulties. This work addresses the challenge by presenting a theoretically grounded alternative for optimizing with relative preference data that circumvents the need for reinforcement learning.

Beyond language-based scenarios, learning policies through preferences has been investigated in both bandit and reinforcement learning domains, leading to a range of proposed methodologies. One notable instance is the contextual dueling bandit (CDB; \cite{yue2012karmed,dudik2015contextual}), which addresses contextual bandit learning by utilizing rankings or pairwise preferences over actions rather than explicit reward feedback. In this framework, where reward signals are unavailable, the concept of optimality is replaced with the \textit{von Neumann winner}—a policy guaranteed to win at least half the time, on expectation, against every competing policy \citep{dudik2015contextual}. However, a distinguishing feature of the CDB paradigm is that preference annotations are provided sequentially in an online manner; in contrast, learning from human preferences typically occurs using a predetermined offline dataset of preference-labeled action pairs \citep{yan2022human}. Likewise, the \textit{preference-based RL} (PbRL) approach leverages binary preference feedback, originating from some unobserved ‘scoring’ function, instead of standard reward signals \citep{BusaFekete2014,ruiz2023dueling}. Several strategies have been introduced for PbRL, some capable of repurposing off-policy preference data; nevertheless, the majority of these methods first estimate the hidden scoring or reward model explicitly before proceeding to policy optimization \citep{jain2013learning,BusaFekete2014,christiano2017deep,sadigh2017active,kupcsik2018learning}. In contrast, we propose a direct, single-stage procedure that learns a policy to fulfill observed preferences without first modeling the reward.

\section{Preliminaries}\label{section:prelims}

We summarize the RLHF process as delineated by \citeauthor{ziegler2020finetuning} and extended in subsequent works \citep{stiennon2022learning, bai2022training, ouyang2022training}. This pipeline generally proceeds in three stages: (1) supervised fine-tuning (SFT); (2) collection of preference data and construction of a reward model; and (3) policy optimization via RL.

\textbf{SFT}: The process typically initiates by further training a pre-trained language model (LM) on a curated dataset, specific to downstream applications such as dialogue or summarization, using supervised learning techniques. This results in the SFT model, denoted by $\pi^\text{SFT}$.

\textbf{Reward Modelling Phase}: During the second stage, prompts $x$ are issued to the SFT model, which generates response pairs $(y_1, y_2)\sim \pi^\text{SFT}(y \mid x)$. Human annotators are then asked to indicate a preference between the responses, recording their choice as $y_w\succ y_l \mid x$, where $y_w$ is the favored and $y_l$ the less preferred output. The presumption is that such preferences reflect an underlying but inaccessible reward model $r^*(y, x)$. Multiple approaches exist for preference modeling; the Bradley-Terry (BT) model \cite{bradley1952rankanalysis} is a prominent example, while the more comprehensive Plackett-Luce model \citep{plackett1975analysis, luce2012individual} may be employed for scenarios involving multiple ranked responses. According to the BT model, the probability distribution over human preferences $p^*$ is formulated as follows:

\begin{equation}\label{eq:bradley-terry}
    p^*(y_1\succ y_2 \mid x)=\frac{\exp\left(r^*(x, y_1)\right)}{\exp\left(r^*(x, y_1)\right) + \exp\left(r^*(x, y_2)\right)}.
\end{equation}

Suppose we are given a fixed dataset of preference comparisons, denoted as $\mathcal{D}=\bigl\{x^{(i)}, y_w^{(i)}, y_l^{(i)}\bigr\}_{i=1}^N$, sampled from the underlying distribution $p^*$. In this case, we can introduce a parameterized reward function $r_{\phi}(x, y)$ and estimate its parameters through maximum likelihood estimation. By recasting the task as a binary classification problem, the objective reduces to minimizing the negative log-likelihood loss:

\begin{equation}\label{eq:reward_model}
    \mathcal{L}_R(r_{\phi}, \mathcal{D}) = -\mathbb{E}_{(x, y_w, y_l)\sim \mathcal{D}}\bigl[\log \sigma(r_{\phi}(x, y_w)- r_{\phi}(x, y_l))\bigr]
\end{equation}

where $\sigma$ denotes the logistic sigmoid function. For language model applications, it is common to initialize the reward network $r_{\phi}(x, y)$ from the parameters of the supervised fine-tuned model $\pi^\text{SFT}(y \mid x)$, appending a linear transformation atop the last hidden state of the transformer, yielding a scalar reward estimate~\cite{ziegler2020finetuning}. To reduce the variance of the reward, prior studies typically normalize the reward outputs such that $\mathbb{E}_{x,y\sim \mathcal{D}}\left[r_\phi(x, y)\right] = 0$ for every $x$.

\textbf{RL Fine-Tuning Phase}: In the reinforcement learning stage, the trained reward function serves as a feedback mechanism for updating the language model. Drawing on previous literature~\citep{jaques2017sequence, jaques2020human}, the training objective can be formulated as:

\begin{equation}\label{eq:RL}
\max_{\pi_{\theta}}  \mathbb{E}_{x\sim \mathcal{D}, y\sim \pi_{\theta}(y \mid x)}\bigl[r_{\phi}(x, y)\bigr] - \beta\mathbb{D}_{\textrm{KL}}\bigl[\pi_{\theta}(y\mid x)\mid \mid \pi_\text{ref}(y\mid x)\bigr],
\end{equation}

with the coefficient $\beta$ governing the allowable divergence from the reference policy $\pi_\text{ref}$—in practice, the pretrained SFT model $\pi^\text{SFT}$. The policy network $\pi_\theta$ is typically initialized from this reference as well. The KL constraint is essential for constraining policy updates, ensuring the resulting policy does not stray too far from regions of the space where the reward model maintains predictive accuracy. Moreover, it discourages mode collapse by preserving generation diversity and mitigates overfitting to spurious high-reward samples. Owing to the nondifferentiable nature of token generation in language models, the objective must be optimized using reinforcement learning algorithms. The prevailing method \citep{ziegler2020finetuning, stiennon2022learning, bai2022training, ouyang2022training} introduces a modified reward function, defined as ${r(x, y) = r_{\phi}(x, y) -\beta (\log \pi_{\theta}(y\mid x) - \log \pi_\text{ref}(y\mid x))}$, and applies Proximal Policy Optimization (PPO)~\cite{schulman2017proximal} for training.

\section{Direct Preference Optimization}\label{sec:DPO}

Seeking to overcome the complexity and computational burden associated with reinforcement learning methods in the context of large-scale language model tuning, we aim to formulate a more streamlined preference-based policy optimization technique. Rather than first estimating a reward and subsequently maximizing it using RL as in traditional RLHF pipelines, we adopt a specialized parameterization for the reward model that admits a closed-form optimal policy, circumventing the need for a full reinforcement learning procedure. As detailed in the following, our principal contribution is to exploit an exact mapping from reward functions to their corresponding optimal policies. This perspective allows us to reparameterize the learning problem, shifting from loss

\textbf{Derivation of the DPO Objective.} Consider the conventional reinforcement learning objective presented in Eq.~\ref{eq:RL}, associated with a general reward function $r$. As detailed in prior research~\citep{peters2007reinforcement, peng2019advantage, korbak2022reinforcement, go2023aligning}, the solution maximizing the expected reward with a KL constraint, as given by Eq.~\ref{eq:RL}, yields the following form:

\begin{equation}\label{eq:op_policy}
    \pi_r(y\mid x) = \frac{1}{Z(x)}\pi_\text{ref}(y\mid x)\exp\left(\frac{1}{\beta}r(x, y)\right),
\end{equation}

Here, $Z(x) =\sum_{y}\pi_\text{ref}(y\mid x)\exp\left(\frac{1}{\beta}r(x, y)\right)$ represents the partition function. A comprehensive derivation is provided in Appendix \ref{app:derivation1}. Notably, substituting even the MLE estimate $r_{\phi}$ for the unknown reward function $r^*$ still leaves the computation of $Z(x)$ intractable for most realistic settings~\citep{korbak2022reinforcement, go2023aligning}, complicating practical application of this formulation. Nevertheless, we can manipulate Eq.~\ref{eq:op_policy} to represent the reward function through the optimal policy $\pi_r$, the reference $\pi_\text{ref}$, and the (potentially uncomputable) partition function $Z(\cdot)$. By taking logarithms of both sides in Eq.~\ref{eq:op_policy} and performing simple algebraic rearrangement, we derive:

\begin{equation}\label{eq:main_eq}
    r(x,y) =\beta \log \frac{\pi_r(y\mid x)}{\pi_\text{ref}(y\mid x)} + \beta \log Z(x).
\end{equation}

This change of variables can be extended to the ground-truth reward $r^*$ as well as the corresponding optimal policy $\pi^*$. Conveniently, the Bradley-Terry model relies only on the relative reward values for two candidate outputs, i.e., ${p^*(y_1 \succ y_2 \mid x) = \sigma(r^*(x, y_1) - r^*(x, y_2))}$. By inserting the alternative form for $r^*(x,y)$ from Eq.~\ref{eq:main_eq} into the preference model given in Eq.~\ref{eq:bradley-terry}, the terms involving the partition function cancel out, leaving a form for the preference probability purely as a function of the optimal policy $\pi^*$ and the reference $\pi_\text{ref}$. Hence, within the Bradley-Terry modeling paradigm, the RLHF-optimal policy $\pi^*$ satisfies the following preference relation:

\begin{equation}\label{eq:objective}
    p^*(y_1\succ y_2 \mid x)=\frac{1}{1 + \exp\left(\beta \log \frac{\pi^*(y_2\mid x)}{\pi_\text{ref}(y_2\mid x)} - \beta \log \frac{\pi^*(y_1\mid x)}{\pi_\text{ref}(y_1\mid x)}\right)}
\end{equation}

Readers can consult Appendix~\ref{app:derivation2} for further details. Although this expression employs the Bradley-Terry preference structure, analogous formulations can be derived under the broader Plackett-Luce models~\citep{plackett1975analysis, luce2012individual}, as demonstrated in Appendix~\ref{app:plackett_luce_models}.

By expressing the preference probabilities using the optimal policy, rather than directly through a reward model, it is now possible to specify a maximum likelihood estimation framework for a parameterized policy $\pi_\theta$. Similar to conventional reward modeling (see Eq.~\ref{eq:reward_model}), the policy-level objective is given by:

\begin{equation}\label{eq:optimum_model}
    \mathcal{L}_\text{DPO}(\pi_{\theta}; \pi_\text{ref}) = -\mathbb{E}_{(x, y_w, y_l)\sim \mathcal{D}}\left[\log \sigma \left(\beta \log \frac{\pi_{\theta}(y_w\mid x)}{\pi_\text{ref}(y_w\mid x)} - \beta \

In this approach, we model an implicit reward via a different parameterization such that the optimal policy aligns directly with $\pi_\theta$. Additionally, because this framework is mathematically equivalent to optimizing a reparametrized Bradley-Terry model, it possesses favorable theoretical guarantees, including consistency given appropriate assumptions about the preference data distribution \cite{bong2022generalized}. Further examination of DPO's theoretical attributes, especially in comparison to related methodologies, is presented in Section~\ref{sec:theory}.

\textbf{Mechanistic perspective on the DPO update.} To better understand how DPO operates, it is insightful to look at the gradient of the loss function $\mathcal{L}_\text{DPO}$. This gradient with respect to the parameter vector $\theta$ is given by:

\begin{multline*}\label{eq:gradient}
    \nabla_\theta \mathcal{L}_\text{DPO}(\pi_\theta;\pi_\text{ref}) = \\ -\beta\mathbb{E}_{(x, y_w, y_l) \sim \mathcal{D}} \bigg[\underbrace{\sigma(\hat{r}_\theta(x, y_l) - \hat{r}_\theta (x, y_w))}_\text{higher weight when reward estimate is wrong}\bigg[\underbrace{\nabla_\theta\log \pi(y_w \mid x)}_\text{increase likelihood of $y_w$} - \underbrace{\nabla_\theta\log\pi(y_l \mid x)}_\text{decrease likelihood of $y_l$}\bigg]\bigg],
\end{multline*}

where $\hat{r}_\theta(x, y) = \beta \log \frac{\pi_\theta(y \mid x)}{\pi_\text{ref}(y \mid x)}$ denotes the implicitly defined reward associated with the language model $\pi_\theta$ relative to the reference $\pi_\text{ref}$ (see Section~\ref{sec:theory} for more details). The interpretation is that the loss gradient amplifies the probability of generating preferred responses $y_w$ while diminishing the probability for less preferred ones $y_l$. Notably, the update is modulated by how much the implicit reward $\hat{r}_\theta$ overvalues dispreferred completions, as measured by $\beta$—essentially reflecting the degree to which the implicit reward model misorders the options and factoring in the KL regularization strength. Empirical results indicate that this adaptive weighting is crucial, as omitting it (i.e., using a straightforward variant that lacks the weighting term) can lead the language model to fail or produce degenerate outputs (refer to Appendix Table~\ref{tab:unlikelihood_generations}).

\textbf{DPO overview.} The typical DPO process consists of two main steps: (1) For each prompt $x$, generate candidate completions $y_1, y_2 \sim \pi_\text{ref}(\cdot \mid x)$ and annotate them with human preference signals to build the offline preference dataset $\mathcal{D} = \{x^{(i)}, y_w^{(i)}, y_l^{(i)}\}_{i=1}^N$; (2) train the language model $\pi_\theta$ by minimizing $\mathcal{L}_\text{DPO}$ with respect to $\pi_\text{ref}$, the dataset $\mathcal{D}$, and a selected $\beta$. In real-world settings, it is desirable to utilize existing public preference datasets instead of collecting new model samples and labels. Because these preference datasets have been sampled from $\pi^\text{SFT}$, we set $\pi_\text{ref} = \pi^\text{SFT}$ whenever possible. In the absence of $\pi^\text{SFT}$, $\pi_\text{ref}$ is initialized by fitting the preferred completions ${(x, y_w)}$ via maximum likelihood, i.e., ${\pi_\text{ref} = \argmax_{\pi}\mathbb{E}_{x, y_w \sim \mathcal{D}}\left[\log \pi(y_w \mid x)\right]}$. This initialization reduces the mismatch between the actual (but unavailable) reference distribution and the $\pi_\text{ref}$ employed by DPO. Implementation specifics and choices of hyperparameters are documented in Appendix~\ref{app:implementation}.

\section{Theoretical Analysis of DPO} This section delivers further conceptual grounding for DPO, establishes its theoretical justification, and situates its strengths relative to actor-critic techniques employed in RLHF (for example, PPO~\cite{schulman2017proximal}).

\label{sec:theory}

\subsection{Your Language Model Is Secretly a Reward Model} DPO forgoes both explicit reward modeling and the use of reinforcement learning in favor of a single maximum likelihood objective. It is important to recognize that the objective function in Eq. \ref{eq:main_eq} coincides with a Bradley-Terry model in which rewards are parameterized by $r^*(x, y) = \beta \log\frac{\pi^*_\theta(y \mid x)}{\pi_\text{ref}(y \mid x)}$. The optimization of the parametric model $\pi_{\theta}$ therefore becomes equivalent, via a change of variables, to the reward model optimization posed in Eq. \ref{eq:reward_model}. In what follows, we present the theoretical framework supporting this reparameterization, demonstrate that it maintains the full representational capacity of reward models, and verify that it can exactly reconstruct the optimal policy. To do so, we start by establishing an equivalence between reward functions.

\begin{definition}
Two reward functions $r(x, y)$ and $r'(x, y)$ are considered equivalent if and only if there exists a function $f$ such that ${r(x, y)-r'(x, y) = f(x)}$.     
\end{definition}

It is straightforward to verify that this indeed constitutes an equivalence relation, thereby dividing the set of reward functions into distinct equivalence classes. This observation allows us to assert the following two lemmas:

\begin{lemma}\label{lemma:same_prefrence}
Within the Plackett-Luce framework, and specifically the Bradley-Terry preference model, any pair of reward functions belonging to the same equivalence class generate identical preference distributions.
\end{lemma}

\begin{lemma}\label{lemma:same_policy}
    Any two reward functions in the same equivalence class yield the same optimal policy when solving the associated constrained RL problem.
\end{lemma}

Since the proofs are elementary, we postpone their presentation to Appendix \ref{app:lemma1}. The first lemma relates to a well-established issue of non-uniqueness in the Plackett-Luce family of preference models \cite{plackett1975analysis}. Because of this inherent ambiguity, it is typical to enforce additional identifiability restrictions to ensure any reliability in MLE estimates from Eq. \ref{eq:reward_model} \cite{bong2022generalized}. The second lemma implies that, for the purpose of policy optimization, it suffices to recover any representative from the relevant equivalence class of reward functions. In Appendix~\ref{app:thm1}, we further demonstrate the following theorem:

\begin{theorem}\label{thm:main}
    Subject to mild regularity conditions, every equivalence class of rewards compatible with the Plackett-Luce (including the Bradley-Terry model) can be expressed using the reparameterization ${r(x, y) = \beta \log \frac{\pi(y\mid x)}{\pi_\text{ref}(y\mid x)}}$ for a particular choice of model $\pi(y\mid x)$ and a given reference policy $\pi_\text{ref}(y \mid x)$.
\end{theorem}

\begin{sproof}
    Take any given reward function $r(x, y)$, which determines an associated optimal policy $\pi_r(y \mid x)$ as per Eq. \ref{eq:op_policy}. We will establish that there exists a member of the equivalence class of $r$ admitting the above reparameterized structure. Specifically, we introduce the projection $f$ by  
\begin{equation}
    f(r; \pi_\text{ref}, \beta)(x, y) = r(x, y) - \beta\log\sum_{y}\pi_\text{ref}(y\mid x)\exp\left(\frac{1}{\beta}r(x, y)\right)
\end{equation}
This mapping $f$ adjusts the original reward by subtracting the log-partition function with respect to $\pi_\text{ref}$. The normalization component depends solely on the context $x$, ensuring that $f(r; \pi_\text{ref}, \beta)(x, y)$ resides within the equivalence class determined by $r(x, y)$. Finally, if we substitute $r$ using the right-hand side of Eq.~\ref{eq:main_eq}, which is valid for any reward function, we find $f(r; \pi_\text{ref}, \beta)(x, y) = \beta \log \frac{\pi_r(y\mid x)}{\pi_\text{ref}(y\mid x)}$. Hence, the projection $f$ generates a member of the original equivalence class exhibiting the target parametric form, implying that no generality is sacrificed in adopting the proposed reparameterization for the reward model.
\end{sproof}

Theorem~\ref{thm:main} can also be interpreted as identifying the specific reward function, among all equivalent formulations, that the DPO reparameterization picks out—namely, the reward function for which

\begin{equation}\label{eq:lag_p}
     \sum_{y}\underbrace{\pi_\text{ref}(y\mid x)\exp\left(\frac{1}{\beta}r(x, y)\right)}_{=\pi(y\mid x)\text{, using Thm.~\ref{thm:main} reparam.}} = 1,
\end{equation}

holds, meaning $\pi(y\mid x)$ forms a legitimate probability distribution (all probabilities are nonnegative and sum to one). If we reference Eq.~\ref{eq:op_policy}, it is evident that Eq.~\ref{eq:lag_p} defines the normalization, or partition function, associated with the optimal policy derived from the reward function $r(x, y)$. The central idea in the DPO algorithm is to enforce constraints on the otherwise ambiguous Plackett-Luce (or more specifically, Bradley-Terry) class of preference models. This constraint maintains the expressive power of possible reward models, while also ensuring that the optimal policy described by Eq.~\ref{eq:op_policy} remains analytically manageable for any prompt $x$.

\subsection{Instability of Actor-Critic Algorithms} Our framework also provides tools for understanding the instability of typical actor-critic methods used in RLHF, such as PPO. Aligning with the RLHF workflow, we specifically examine the reinforcement learning fine-tuning phase discussed in Section \ref{section:prelims}. Our approach draws upon the control as inference paradigm \cite{levine2018reinforcement} within the context of the constrained RL setup in \ref{eq:RL}. We suppose a family of parameterized models $\pi_{\theta}(y\mid x)$ and aim to minimize $\mathbb{D}_{\text{KL}}[\pi_{\theta}(y|x) \mid \mid \pi^*(y\mid x)]$, where $\pi^*$ denotes the optimal policy described by Eq.~\ref{eq:optimum_model}, determined by the reward $r_{\phi}(y, x)$. This leads, after some derivation, to the following optimization objective:

\begin{equation}\label{eq:AC}
    \max_{\pi_{\theta}}\mathbb{E}_{\pi_{\theta}(y\mid x)}\bigg[\underbrace{r_{\phi}(x, y) -\beta\log\sum_{y}\pi_\text{ref}(y\mid x)\exp\left(\frac{1}{\beta}r_{\phi}(x, y)\right)}_{f(r_{\phi}, \pi_\text{ref}, \beta)} - \underbrace{\beta\log\frac{\pi_{\theta}(y\mid x)}{\pi_\text{ref}(y\mid x)}}_{\text{KL}}\bigg]
\end{equation}

The objective considered here is identical to that explored in previous research 
\citep{ziegler2020finetuning, stiennon2022learning, bai2022training, ouyang2022training}, wherein the DPO-equivalent reward is applied to the reward class $r_{\phi}$. In this context, the normalization component in $f(r_{\phi}, \pi_\text{ref}, \beta)$ can be viewed as the soft value function associated with the reference policy $\pi_\text{ref}$. Although this component does not alter the set of optimal solutions, omitting it can result in policy gradients with elevated variance, thereby destabilizing the learning process. One potential remedy is to utilize a learned value function for normalization, yet this approach may also introduce optimization challenges. Alternatively, earlier approaches have used human completions as baselines, providing a single-sample Monte Carlo approximation to the normalization term. In comparison, the DPO reparameterization produces a reward formulation that eliminates the need for such baselines.

\section{Experiments}
Here, we present an empirical assessment of DPO’s effectiveness in training policies directly from preference data. We start by investigating a controlled text-generation environment to answer the following: how well does DPO balance reward maximization against minimizing KL-divergence from the reference policy, relative to established preference learning methods such as PPO? We subsequently apply DPO to larger-scale models and more complex RLHF scenarios, including tasks like summarization and dialogue. Our findings indicate that DPO typically matches or surpasses competitive baselines, such as RLHF implemented with PPO and strategies that select the best out of $N$ sampled trajectories under a trained reward model, all with minimal hyperparameter adjustment. The subsequent section details the experimental protocols; further specifics can be found in Appendix~\ref{app:exp_details}.

\textbf{Tasks.} We conduct experiments on three distinct open-ended text generation tasks. Across all tasks, methods learn a policy from a set of preference-labeled data, $\mathcal{D}=\bigl\{x^{(i)}, y_w^{(i)}, y_l^{(i)}\bigr\}_{i=1}^N$. In the \textbf{controlled sentiment generation} task, $x$ serves as the opening of a movie review from the IMDb corpus \cite{maas-EtAl:2011:ACL-HLT2011}, with the objective that the policy generates a completion $y$ that exhibits positive sentiment. To enable systematic evaluation, we generate preference pairs over completions using an existing sentiment classifier, ensuring that $p(\text{positive}\mid x,y_w)>p(\text{positive}\mid x,y_l)$. For SFT, GPT-2-large is fine-tuned to convergence on training-set reviews from IMDb (details are provided in App~\ref{app:sentiment_details}). In the case of \textbf{summarization}, $x$ consists of a Reddit forum post, and the policy's goal is to generate a concise summary $y$ highlighting the post’s key points. Adopting previous approaches, we employ the Reddit TL;DR summarization dataset \citep{volske-etal-2017-tl} augmented with human preference annotations from \citeauthor{stiennon2022learning}. For SFT, we utilize a model adapted to human-created summaries of forum posts,\footnote{\url{https://huggingface.co/CarperAI/openai_summarize_tldr_sft}} using the TRLX RLHF framework \citep{leandro_von_werra_2023_7790115}. The preference annotations come from \citeauthor{stiennon2022learning}, who collected them on samples generated by a comparable, but separately trained, SFT model. Lastly, in the \textbf{single-turn dialogue} setting, $x$ represents a user query spanning a wide range of topics from scientific inquiries to personal advice. Here, the policy must deliver a relevant and engaging response $y$; we employ the Anthropic Helpful and Harmless dialogue dataset \citep{bai2022training}, which features 170k human–assistant dialogues. Each dialogue concludes with two model-generated responses and an accompanying label indicating which one a human annotator preferred. Because there is no available SFT checkpoint in this scenario, we fine-tune a general-purpose language model exclusively on preferred responses to construct the SFT baseline.

\textbf{Evaluation.} We employ two distinct evaluation strategies in our experiments. To rigorously assess each algorithm's ability to optimize the constrained reward maximization objective, we utilize a controlled sentiment generation environment. Here, we examine the trade-off frontier between the achieved reward and KL-divergence from a reference policy, leveraging the availability of a known, ground-truth reward function provided by a sentiment classifier. Conversely, since the true reward function is typically inaccessible in practical scenarios, we rely on \textit{win rate} comparisons against a baseline policy for real-world tasks. In this context, GPT-4 serves as a stand-in for human evaluation, measuring summary quality for summarization tasks and response helpfulness in single-turn dialogue settings. For the summarization domain, the comparison baseline is the set of reference summaries from the test split; in dialogue, we benchmark against the preferred human-labeled response from the test set. While recent research has indicated that LMs can serve as superior automated evaluators compared to standard metrics \citep{Chen2023ExploringTU}, we also undertake a human subject study to substantiate our decision to use GPT-4, as detailed in Sec.~\ref{sec:human-judgments}. Our findings reveal a high correlation between GPT-4 and human judgments, with human agreement with GPT-4 as strong as, or stronger than, inter-annotator consistency among humans.

\textbf{Methods.} Beyond DPO, we assess multiple prevalent strategies for training language models to align with human preferences. For the summarization task, we apply zero-shot prompting to \textbf{GPT-J} \citep{gpt-j}, whereas, for dialogue, we employ 2-shot prompting using \textbf{Pythia-2.8B} \citep{biderman2023pythia}. We further evaluate the \textbf{SFT} approach as well as \textbf{Preferred-FT}, which involves supervised fine-tuning on the selected completion $y_w$ derived either from the SFT model (for controlled sentiment and summarization) or from a generic LM (in single-turn dialogue scenarios). Another variant based on pseudo-supervised learning is the \textbf{Unlikelihood} method~\citep{welleck2019neural}, which modifies the objective to boost the likelihood of $y_w$ while simultaneously suppressing the likelihood of $y_l$, regulated by an optional `unlikelihood' term coefficient $\alpha\in[0,1]$. We also incorporate \textbf{PPO} \citep{schulman2017proximal}, which uses a reward model inferred from preference data, and an oracle version, \textbf{PPO-GT}, which directly leverages the true reward function present in the controlled sentiment scenario. For PPO-GT in sentiment-related experiments, we test both a standard implementation \cite{leandro_von_werra_2023_7790115} and a tailored variant featuring normalized rewards and refined hyperparameters for enhanced results; these modifications are equally applied to PPO runs with learned reward functions. Lastly, we include a \textbf{Best of $N$} strategy, wherein $N$ candidate responses are sampled from the SFT model (or from Preferred-FT for dialogue), and the response with the highest reward score, as determined by a learned reward model, is selected. While this method achieves strong empirical results and sidesteps the dependency of reward quality on PPO optimization, it is generally infeasible for larger values of $N$ because it demands the sampling of $N$ outputs for each test input.

\subsection{How well can DPO optimize the RLHF objective?}

In standard RLHF approaches, the objective function maximizes the reward while enforcing a KL constraint to prevent the policy from diverging significantly from the reference policy. Thus, comparing different methods necessitates jointly evaluating both the reward and the associated KL divergence; maximizing reward at the cost of a substantial increase in KL is not inherently preferable. Figure~\ref{fig:frontier-tldr-main} illustrates the trade-off between reward and KL across several algorithms within the sentiment analysis context. For each algorithm, we conduct several independent runs, each with a distinct setting of the conservativeness hyperparameter (PPO uses target KL values $\in\{3,6,9,12\}$, unlikelihood is run with $\beta \in \{0.05,0.1,1,5\}$, and $\alpha\in\{0.05,0.1,0.5,1\}$; preferred-FT is varied by random seed). This experimental sweep consists of 22 runs overall. At intervals of 100 training steps and continuing until training converges, each resulting policy is evaluated on a test prompt set by computing the mean reward (according to the true reward function) and the mean sequence-level KL\footnote{Specifically, this is the aggregate of KL-divergences at each timestep.} against the reference policy, i.e., $\text{KL}\left(\pi\mid \mid \pi_\text{ref}\right)$. The results demonstrate that DPO attains the most favorable frontier, simultaneously reaching higher rewards and maintaining low KL. Notably, this outcome stands out for several reasons. Despite both DPO and PPO optimizing an identical objective, DPO significantly outperforms PPO in terms of efficiency; DPO’s trade-off between reward and KL unambiguously surpasses that of PPO. Moreover, DPO continues to deliver a superior reward-KL frontier even in comparison to PPO variants that utilize access to ground truth rewards (PPO-GT).

\subsection{Can DPO scale to real preference datasets?}
\label{sec:dpo-real-datasets}
We proceed by assessing the fine-tuning effectiveness of DPO on tasks involving summarization and single-turn dialogue. In the context of summarization, it is well documented that common automatic metrics like ROUGE often show weak alignment with human judgements~\citep{stiennon2022learning}. Prior research has demonstrated that training language models with PPO guided by human preferences yields more suitable summaries. To compare various approaches, we generate completions for examples in the TL;DR summarization dataset's test set and calculate the average win rate relative to reference completions. Samples across all approaches are drawn using a temperature range between 0.0 and 1.0; corresponding win rates are displayed in Figure~\ref{fig:frontier-tldr-main} (right). DPO, PPO, and Preferred-FT each continue training from the same GPT-J SFT baseline model\footnote{\url{https://huggingface.co/CarperAI/openai_summarize_tldr_sft}}. Our analysis shows that DPO achieves a win rate of roughly 61\% at temperature 0.0, surpassing PPO, which reaches its peak at about 57\% at the same temperature. In addition, DPO attains a higher best-case win rate than the strongest $N$ baseline. We emphasize that minimal effort was made to tune DPO's $\beta$ hyperparameter, so the reported outcomes could be an underrepresentation of DPO's true capabilities. Furthermore, DPO exhibits greater resilience to variations in sampling temperature than PPO, whose win rates can fall to match that of the original GPT-J at elevated temperatures. The Preferred-FT approach yields only marginal improvements over the SFT baseline. A direct human evaluation comparison between DPO and PPO is presented in Section~\ref{sec:human-judgments}, where samples from DPO at temperature 0.25 were favored 58\% of the time relative to PPO outputs at temperature 0.

For single-turn dialogue evaluation, we assess various approaches on the subset of the Anthropic HH dataset's test split \citep{bai2022training} containing only one round of human-assistant exchange. To measure method win rates, GPT-4 judges responses using the reference of preferred completions from the test set. Due to the absence of a canonical SFT model for this scenario, we initiate our experiments with a pre-trained Pythia-2.8B and employ Preferred-FT to train a reference model using the selected completions, ensuring the outputs align with the model's distribution. Subsequent DPO training is conducted on top of this setup. For comparison, we consider the best completion among 128 samples from Preferred-FT (noting that performance improvements plateau at 128 samples for this benchmark; consult Appendix Figure~\ref{fig:best-of-n}) as well as a 2-shot prompted variant of the base Pythia-2.8B. Results indicate that DPO either matches or surpasses these baselines at optimal temperatures for each method. Additionally, we test an RLHF model, trained via PPO on the Anthropic HH dataset \footnote{\url{https://huggingface.co/reciprocate/ppo_hh_pythia-6B}}, provided by a reputable repository \footnote{\url{https://github.com/CarperAI/trlx/tree/main/examples/hh}}. However, we are unable to identify any prompt or sampling temperature that allows it to outperform the original Pythia-2.8B model. Taking into account findings from TL;DR and the alignment in reward functions, we use Best of 128 as an approximate indicator of PPO-level outcomes. In summary, DPO emerges as the sole method with both computational efficiency and consistent improvement over preferred completions within the Anthropic HH dataset, performing on par with or better than the resource-intensive Best of 128 baseline. Lastly, Figure~\ref{fig:dialogue-main} illustrates DPO's rapid convergence to peak performance.

\subsection{Generalization to a new input distribution}

To further assess the robustness of PPO and DPO under distributional changes, we evaluate the policies obtained from our Reddit TL;DR summarization experiments on an out-of-distribution dataset: news articles from the test split of CNN/DailyMail \citep{nallapati-etal-2016-abstractive}, applying the optimal sampling temperatures from TL;DR (0 and 0.25). Table~\ref{tab:ood} summarizes the findings. We measured the rate at which GPT-4 selected system summaries over reference summaries using the identical GPT-4 (C) prompt as in the Reddit TL;DR evaluations, substituting ``forum post'' with ``news article''. On this alternate dataset, DPO maintains a clear performance advantage over PPO. These results provide preliminary evidence that DPO-trained policies can generalize at least as well as PPO-trained ones, despite the absence of the additional unlabeled Reddit TL;DR prompts available to PPO during training.

\subsection{Validating GPT-4 judgments with human judgments}
\label{sec:human-judgments}
To assess the trustworthiness of GPT-4 as an evaluator, we ran a human study leveraging summaries from the TL;DR experiment and two distinct GPT-4 prompt styles. The \textbf{GPT-4 (S)} (simple) prompt requests a judgment of which summary best encapsulates the post’s main points. By contrast, the \textbf{GPT-4 (C)} (concise) prompt asks for both informativeness and conciseness; we include this prompt due to an observed GPT-4 bias toward generating longer, more redundant summaries under the \textbf{GPT-4 (S)} prompt relative to human raters. Appendix~\ref{app:prompts} contains full prompt texts. For diversity in sample quality, we evaluate three conditions: the strongest (DPO, temp. 0.25), the weakest (PPO, temp. 1.0), and a middle-performing (SFT, temp. 0.25) method, each compared to PPO samples at the most effective temperature (greedy sampling). Across both prompts, GPT-4’s preferences typically align with those of human annotators to a similar degree as humans agree among themselves, indicating GPT-4’s reliability as a human stand-in (with repeated human ratings gathered for DPO and PPO-1 only due to annotator limitations). Notably, the \textbf{GPT-4 (C)} prompt yields agreement and win rates most reflective of human judgments; therefore, we report results for Section~\ref{sec:dpo-real-datasets} using this prompt. Detailed information about the human study procedure, including the interface and the pool of annotators, appears in Appendix~\ref{app:human-study}.

\section{Discussion}
Learning from preferences constitutes an effective and highly scalable methodology for developing competent and value-aligned language models. In this work, we have presented DPO, an uncomplicated training approach that leverages preferences to guide language model development without resorting to reinforcement learning techniques. Unlike existing methods that transform the preference learning challenge into a conventional RL framework to utilize standard RL algorithms, DPO establishes a correspondence between language model policies and reward representations, thereby permitting direct training of models to adhere to human preferences through a straightforward cross-entropy loss—eliminating the need for reinforcement learning and maintaining generality. Remarkably, with little to no adjustment of hyperparameters, DPO achieves results that are on par with or superior to leading RLHF strategies, such as those based on PPO, thereby substantially lowering the barriers associated with preference-based language model training.

\textbf{Limitations \& Future Work.} Our findings prompt several directions for further investigation. One key question is the extent to which DPO-trained policies demonstrate out-of-distribution generalization, particularly in comparison to models optimized with explicit reward functions. Preliminary evidence indicates DPO models may generalize comparably to their PPO-trained counterparts; however, additional systematic analysis is warranted. For instance, it remains to be explored whether self-labeling strategies based on DPO can harness unlabeled prompts with similar efficacy. Another line of inquiry concerns the dynamics of reward over-optimization under direct preference optimization, and whether the modest drop in performance illustrated in Figure~\ref{fig:dialogue-main}-right exemplifies this phenomenon. Although our current evaluations consider models up to 6B parameters, scaling DPO to much larger, cutting-edge models represents a promising avenue for future research. In terms of assessment protocols, we observe that GPT-4-based win rates are sensitive to prompting; subsequent studies may investigate methods for eliciting more reliable evaluations from automated systems. Lastly, DPO holds promise for broader applications beyond language modeling from human preferences, such as training generative models across different domains. 

\end{document}